\chapter{Limite centrale, orbites de Kepler et holonomie de Sommerfeld}
\label{chap:limite-central-kepler-sommerfeld}
\index{Kepler, orbite de}\index{Sommerfeld, quantification de}
\index{Levi-Civita, revêtement de}\index{Möbius, holonomie de}

Ce chapitre distingue quatre strates que la terminologie physique peut
superposer indûment : l'annulation exacte des flux intérieurs dans un
graphe ; la réalisation continue d'un champ central ; la dynamique képlérienne
qui résulte du potentiel inverse ; et la condition semi-classique de phase sur
une bande avec holonomie. La première est combinatoire. Les trois autres
exigent successivement une carte métrique, un lagrangien et une ligne de
phase semi-classique.

\section{Gauss discret et descente inverse à l'aire}
\label{sec:gauss-discreto-limite-central}

Soit \(G=(V,E)\) un graphe fini non orienté, soit
\(\vec E:=\{(x,y),(y,x):\{x,y\}\in E\}\) son ensemble d'arcs et soit
\(F:\vec E\to\mathbb R\) une \(1\)-co\-chaîne antisymétrique,
\(F(y,x)=-F(x,y)\). Nous définissons
\[
 \operatorname{div}F(x)
 :=\sum_{y:\{x,y\}\in E}F(x,y).
\]
Pour \(S\subseteq V\), le flux
\(\operatorname{Flux}_{\partial S}F\) est la somme orientée sur les arêtes
ayant une extrémité dans \(S\) et l'autre dans \(V\setminus S\).

\begin{theorem}[Théorème de Gauss sur les graphes]
\label{thm:gauss-grafo-hmt-md}
Pour tout \(S\subseteq V\),
\[
 \sum_{x\in S}\operatorname{div}F(x)
 =\operatorname{Flux}_{\partial S}F.
 \tag{K1}\label{eq:gauss-grafo-hmt-md}
\]
\end{theorem}

\begin{proof}
Chaque arête dont les deux extrémités sont dans \(S\) apparaît deux fois avec des orientations
opposées et s'annule. Seules subsistent les arêtes qui traversent la chaîne
de bord.
\end{proof}

Dans le réseau cubique, pour \(R\in\mathbb Z_{\geq0}\),
\[
 C_R:=\{x\in\mathbb Z^3:\|x\|_\infty\leq R\},
\]
le nombre de liens sortants est
\[
 |\partial C_R|=6(2R+1)^2.
\]
Si une source possède un flux total \(Q\) et que l'on impose une uniformité moyenne sur la
coquille, le flux par lien vaut
\[
 E_R=\frac{Q}{6(2R+1)^2}.
\]
En définissant \(r_{\rm eff}\) par
\[
 4\pi r_{\rm eff}^2=6(2R+1)^2,
\]
on obtient l'identité exacte
\[
 E_R=\frac{Q}{4\pi r_{\rm eff}^2}.
 \tag{K2}\label{eq:inverso-area-cascaron}
\]
Le passage de \eqref{eq:gauss-grafo-hmt-md} à
\eqref{eq:inverso-area-cascaron} utilise l'hypothèse supplémentaire d'isotropie
moyenne. Identifier \(Q\) à une source gravitationnelle ou électrique exige,
en outre, le transducteur dimensionnel correspondant.

\section{Lagrangien central et coniques de Kepler}
\label{sec:lagrangiano-central-kepler}

Soient \(\mu>0\) une masse réduite et
\(\mathcal K>0\) un couplage de type produit d'énergie et de longueur,
\([\mathcal K]=\mathsf E\,\mathsf L\). Dans le feuillet
euclidien continu, considérons
\[
 L_{\rm Kep}(\mathbf r,\dot{\mathbf r})
 =\frac{\mu}{2}\|\dot{\mathbf r}\|^2+\frac{\mathcal K}{r},
 \qquad
 H_{\rm Kep}(\mathbf r,\mathbf p)
 =\frac{\|\mathbf p\|^2}{2\mu}-\frac{\mathcal K}{r},
 \tag{K3}\label{eq:lagrangiano-hamiltoniano-kepler}
\]
où \(r=\|\mathbf r\|>0\) et
\(\mathbf p=\mu\dot{\mathbf r}\). L'équation d'Euler--Lagrange est
\[
 \mu\ddot{\mathbf r}
 =-\frac{\mathcal K}{r^3}\mathbf r.
 \tag{K4}\label{eq:ecuacion-central-kepler}
\]
Le moment cinétique spécifique
\(\mathbf h=\mathbf r\times\dot{\mathbf r}\) est constant ; si
\(\mathbf h\ne0\), la trajectoire reste dans le plan orthogonal à
\(\mathbf h\).

\begin{proposition}[Conique képlérienne]
\label{prop:conica-kepleriana-hmt-md}
Soit \(h=\|\mathbf h\|\). Toute solution non radiale de
\eqref{eq:ecuacion-central-kepler} satisfait
\[
 r(\phi)
 =\frac{p}{1+e\cos(\phi-\phi_0)},
 \qquad
 p=\frac{\mu h^2}{\mathcal K}.
 \tag{K5}\label{eq:conica-kepleriana}
\]
Pour une énergie négative, \(0\leq e<1\) et l'orbite est elliptique.
\end{proposition}

\begin{proof}
Avec \(u=1/r\), l'équation de Binet prend la forme
\[
 u''(\phi)+u(\phi)=\frac{\mathcal K}{\mu h^2}.
\]
Sa solution générale est
\[
 u=\frac{\mathcal K}{\mu h^2}
 \bigl(1+e\cos(\phi-\phi_0)\bigr).
\]
Inverser \(u\) donne \eqref{eq:conica-kepleriana} ; la condition de liaison
équivaut à \(e<1\).
\end{proof}

La vitesse aréolaire est
\[
 \frac{\mathrm dA}{\mathrm dt}=\frac h2.
 \tag{K6}\label{eq:segunda-ley-kepler}
\]
Pour une ellipse de demi-grand axe \(a\), intégrer
\eqref{eq:segunda-ley-kepler} sur une période produit
\[
 A_{\rm ell}=\pi ab,
 \qquad
 b=a\sqrt{1-e^2},
 \qquad
 p=a(1-e^2)=\frac{\mu h^2}{\mathcal K}.
\]
Par conséquent,
\(T=2A_{\rm ell}/h\), et en éliminant \(b\), \(p\) et \(h\), on obtient
\[
 T^2=\frac{4\pi^2\mu}{\mathcal K}\,a^3.
 \tag{K7}\label{eq:tercera-ley-kepler}
\]
Dans le problème gravitationnel d'épreuve,
\(\mathcal K=G M\mu\), et
\eqref{eq:tercera-ley-kepler} se réduit à
\(T^2=4\pi^2a^3/(GM)\) \cite{Kepler1609,Newton1687}.

Le théorème de Gauss du graphe est exact. Sous les isotropie et distance
déclarées, le flux produit le champ inverse à l'aire ; la limite continue
transporte ensuite ce champ vers le lagrangien central et les orbites
képlériennes écrites ci-dessus.

\section{Revêtement de Levi--Civita et ruban de Möbius}
\label{sec:levi-civita-mobius-kepler}

Dans le plan complexe privé de l'origine, l'application de Levi--Civita
\[
 q=z^2
 \tag{K8}\label{eq:cubierta-levi-civita}
\]
est un revêtement double. Un tour de \(q\) autour de l'origine échange
les deux relèvements \(z\) et \(-z\) ; deux tours sont nécessaires pour fermer le
relèvement. C'est
un témoin exact de mémoire \(\mathbb Z_2\). Ce n'est cependant pas un ruban
de Möbius : la régularisation complète du problème de Kepler exige aussi
de reparamétrer le temps et de transformer les moments.

Soit maintenant \(\gamma:S^1\hookrightarrow\mathbb R^2\) un plongement dont l'image
est une orbite elliptique, et fixons \(\varepsilon>0\). Nous définissons
\[
 M_\gamma
 :=([0,2\pi]\times[-\varepsilon,\varepsilon])/
 \bigl((0,u)\sim(2\pi,-u)\bigr).
 \tag{K9}\label{eq:banda-mobius-orbita}
\]

\begin{theorem}[Période orientée du ruban orbital]
\label{thm:periodo-orientado-banda-orbital}
\(M_\gamma\) est un ruban de Möbius dont le cœur est \(\gamma\). Une section
de la ligne de signe possède l'holonomie
\[
 \mathcal H(\gamma)=-I,
 \qquad
 \mathcal H(\gamma^2)=I.
 \tag{K10}\label{eq:holonomia-banda-orbital}
\]
Par conséquent, son retour comme section orientée exige deux tours.
\end{theorem}

\begin{proof}
L'identification des fibres extrêmes par \(u\mapsto-u\) est la
construction standard du ruban de Möbius. L'holonomie d'un tour est
l'involution de la fibre ; son carré est l'identité.
\end{proof}

Le revêtement \eqref{eq:cubierta-levi-civita} et le ruban
\eqref{eq:banda-mobius-orbita} partagent une monodromie d'ordre deux, mais ne
sont pas des objets isomorphes. Cette distinction empêche de déduire une orientation
particule--antiparticule de la régularisation de Levi--Civita sans une
application supplémentaire.

\section{Condition semi-classique de Sommerfeld avec holonomie}
\label{sec:sommerfeld-holonomia-banda}

Soit
\[
 J_\gamma:=\oint_\gamma\mathbf p\cdot\mathrm d\mathbf q
\]
l'action d'un cycle et supposons que le transport semi-classique apporte la
phase \(\exp(\mathrm iJ_\gamma/\hbar_{\rm int})\). Si la ligne réelle associée
possède l'holonomie \(\varepsilon_\gamma\in\{+1,-1\}\), la condition de clôture de
la section complexifiée est
\[
 \varepsilon_\gamma
 \exp\!\left(\frac{\mathrm iJ_\gamma}{\hbar_{\rm int}}\right)=1.
 \tag{K11}\label{eq:cierre-fase-holonomia-sommerfeld}
\]

\begin{proposition}[Décalage d'une demi-unité par holonomie de signe]
\label{prop:desplazamiento-sommerfeld-holonomia}
Sous les hypothèses précédentes,
\[
 \begin{array}{lll}
 \varepsilon_\gamma=+1
 &\Longrightarrow&
 J_\gamma=2\pi\hbar_{\rm int}n,\\[1mm]
 \varepsilon_\gamma=-1
 &\Longrightarrow&
 J_\gamma=2\pi\hbar_{\rm int}\left(n+\dfrac12\right),
 \end{array}
 \qquad n\in\mathbb Z.
 \tag{K12}\label{eq:cuantizacion-sommerfeld-holonomia}
\]
\end{proposition}

\begin{proof}
La \cref{eq:cierre-fase-holonomia-sommerfeld} exige une phase égale à un dans
le secteur trivial et égale à moins un dans le secteur de signe. Leurs arguments
diffèrent, respectivement, d'un multiple entier ou demi-entier de \(2\pi\).
\end{proof}

La proposition est une conséquence exacte de la phase et de l'holonomie
déclarées ; elle ne sélectionne pas à elle seule le Hamiltonien de Coulomb, un nombre
quantique physique ni un spectre atomique \cite{Sommerfeld1916}. Dans une
quantification EBK, il faut ajouter les indices de Maslov des cycles
correspondants, sans compter deux fois le décalage topologique.

Relativement au lecteur d'orientation, les deux feuillets locaux peuvent
représenter les chemins \(\gamma\) et \(\gamma^\vee\) et se relier aux
deux orientations de l'intrication bicouche définie dans
\cref{def:entropia-bicapa-rev6}. Cette correspondance n'identifie pas
ces chemins à des fonctions de Green avancées ou retardées et n'affirme pas que
l'entropie cause la conique képlérienne : l'orbite provient de
\eqref{eq:ecuacion-central-kepler} ; l'entropie mesure, après une
complétion de Hilbert, la corrélation entre feuillets.

\begin{criterion}[Limite Kepler--Sommerfeld]
\label{crit:limite-kepler-sommerfeld}
La réalisation est réfutée ou indûment typée si l'on obtient une loi
en inverse du carré sans déclarer l'isotropie ; si
\(\mathcal K\) ou \(G\) sont introduits comme scalaires sans droite dimensionnelle ;
si le revêtement \(q=z^2\) est identifié à un ruban de Möbius ; si
\eqref{eq:cuantizacion-sommerfeld-holonomia} est publié comme spectre exact
sans Hamiltonien, cycles et indices de Maslov ; ou si l'on attribue
une intrication à une simple inversion déterministe de chemin.
\end{criterion}
